\section{Finite mixed polynomials and buffered pinning}\label{poly:section}

We develop a rule for fixing one vertex's color while controlling the
total positive-root mass of an expected characteristic polynomial.
The rule will be applied repeatedly in the measured paving argument;
its constants must therefore be independent of the matrix size and
of previous color choices. We first identify the rooted mass and the
change caused by one pin, then prove the inequality that relates them.

\subsection{Stable polynomials and local spectral data}

Let $Q$ be a finite set, let $W=(W_{vw})_{v,w\in Q}$ be Hermitian
with zero diagonal, and let $a_v\in\R$. At each $v$, independently choose
a color in $\{1,\ldots,r\}$ with probabilities $t_{vi}\geq0$,
$\sum_i t_{vi}=1$. The matrix $W_{\mathrm{pav}}$ retains the entries between
vertices of the same color and sets all other entries to zero. Define
\begin{equation}\label{poly:definition}
 f_Q(Y)=\E\det\bigl(\diag(Y_v+a_v)_{v\in Q}-W_{\mathrm{pav}}\bigr),
 \qquad f_\varnothing=1.
\end{equation}
All restrictions to subsets of $Q$ use the restricted matrix, penalties,
and probabilities. The polynomial is real and multiaffine, its leading
monomial is $\prod_{v\in Q}Y_v$, and
\begin{equation}\label{poly:deletion}
 \partial_{Y_v}f_Q=f_{Q\setminus\{v\}}.
\end{equation}
We write $f_Q(z)$ when every variable is set equal to $z$.
The uniform-probability paving polynomial and its determinant-derivative
representation occur in \cite[Theorem~1.6 and Proposition~1.7]{RL20};
independent nonuniform distributions are also discussed in
\cite[remark following Lemma~24]{RS21}. We give the identities and
stability arguments in the form needed for local perturbations below.

\begin{lemma}[Stability]\label{lem:stability}
Put $t_i^{1/2}=\diag(\sqrt{t_{vi}})_{v\in Q}$. Then
\begin{equation}\label{poly:differential-formula}
 f_Q(Y)=
 \left(\prod_{v\in Q}\frac{1}{r!}\partial_{Y_v}^{r-1}\right)
 \prod_{i=1}^r
 \det\bigl(\diag(Y+a)-r t_i^{1/2}Wt_i^{1/2}\bigr).
\end{equation}
This identity remains valid for complex $t_{vi}$ with $\sum_i t_{vi}=1$,
where the expectation means the finite algebraic sum with the corresponding
product weights, and the same square root is used in both occurrences of
$t_i^{1/2}$. The algebraic identity also holds with complex penalties
and arbitrary complex off-diagonal entries of $W$, still with zero diagonal.

More generally, let $\mathcal H$ be an open half-plane. If the product of
determinants on the right of \eqref{poly:differential-formula}, for every
principal restriction, is nonzero whenever all its variables lie in
$\mathcal H$, then the same holds for $f_Q$. In particular, for Hermitian
$W$ and real probabilities and penalties, $f_Q$ is upper-half-plane stable and
its univariate specialization is real-rooted.
\end{lemma}

\begin{proof}
Expand each determinant by permutations. Since $W$ has zero diagonal,
the off-diagonal terms are collections of disjoint nontrivial permutation
cycles. Each such cycle $C$ has sign minus and weight
\begin{equation}\label{poly:cycle-weight}
 w_C=\left(\prod_{(v,w)\in C}W_{vw}\right)
 \sum_{i=1}^r\prod_{v\in C}t_{vi}
\end{equation}
in the expectation; unused vertices contribute $Y_v+a_v$. Indeed, all
vertices on a cycle must receive one common color, and disjoint cycles
involve disjoint independent choices.

In the product on the right of \eqref{poly:differential-formula}, suppose
a vertex occurs in nontrivial cycles of $k$ different determinant factors.
Its degree in $Y_v+a_v$ is $r-k$. If $k\geq2$, the indicated derivative
kills the term. If $k=0$, the normalized derivative sends
$(Y_v+a_v)^r$ to $Y_v+a_v$. If $k=1$, it sends
$(Y_v+a_v)^{r-1}$ to $1/r$, cancelling the factor $r$ from the matrix
scaling at that vertex. The remaining terms are precisely
\eqref{poly:cycle-weight} for disjoint cycles. The two square-root factors
at each cycle vertex multiply to $t_{vi}$, so this argument also proves
the complex-parameter identity independently of the choices of square roots.
This cycle calculation is algebraic in the entries of $W$ and does not
require them to be Hermitian; that hypothesis is used for the real-data
stability conclusion, not for the identity.

For the half-plane assertion, proceed by induction on the subset $D$ of
vertices at which $r-1$ derivatives have already been taken, simultaneously
for all principal restrictions. Fix the other variables in $\mathcal H$
and take a vertex $v\notin D$. In each determinant, the coefficient of
$Y_v$ is its principal determinant with $v$ removed, even for complex
penalties. Extracting the coefficient of $Y_v^r$ from the product
commutes with all normalized derivatives at vertices in $D$. Thus this
coefficient is exactly the already-processed product for the principal
set with $v$ removed. It is
nonzero by induction, so the polynomial in $Y_v$ has degree exactly $r$.
Its roots lie in the closed convex complement of $\mathcal H$.
Gauss--Lucas implies that each of its successive derivatives has all roots
in that complement and is nonzero. This proves the induction step.

For completeness, the determinant hypothesis holds whenever
$\mathcal H$ lies strictly beyond a numerical-range bound for
\begin{equation}\label{poly:matrix-numerical-range}
 B_i=r t_i^{1/2}Wt_i^{1/2}-\diag(a).
\end{equation}
For example, if $\operatorname{Im}\ip{B_i u}{u}\leq\beta\norm{u}^2$
and $\operatorname{Im}Y_v>\beta$ for all $v$, a null vector of
$\diag(Y)-B_i$ contradicts the imaginary part of its quadratic form.
The same reasoning applies after a common rotation of the half-plane;
principal compressions inherit numerical-range bounds. For Hermitian
$W$ and real probabilities and penalties, the matrices
\eqref{poly:matrix-numerical-range} are
Hermitian, giving upper-half-plane stability. Real coefficients and
complex conjugation then give real-rootedness of the univariate
specialization.
\end{proof}

\begin{lemma}[Local spectral data]\label{lem:local-data}
Fix $v\in Q$, put $f=f_Q$, $g=f_{Q\setminus\{v\}}$, and write
$h=f/g$, $G=1/h$ at univariate arguments. For $z\in\C$ with
$\operatorname{Im}z>0$,
\begin{equation}\label{poly:pick-inequality}
 \operatorname{Im}h(z)\geq\operatorname{Im}z.
\end{equation}
If $s_j$ are the roots of $g$, with multiplicities, there are
$\nu_j\geq0$ such that
\begin{equation}\label{poly:arrow-quotient}
 h(z)=z+a_v-\sum_j\frac{\nu_j}{z-s_j}.
\end{equation}
Consequently, on choosing $|c_j|^2=\nu_j$,
\begin{equation}\label{poly:arrow-matrix}
 B=\begin{pmatrix}-a_v&c^*\\ c&S\end{pmatrix},
 \qquad S=\diag(s_j),
 \qquad \det(zI-B)=f(z).
\end{equation}
The function $G$ is the transform
\begin{equation}\label{poly:stieltjes}
 G(z)=\int_\R\frac{d\gamma_v(x)}{z-x}
\end{equation}
of the spectral probability measure of $B$ at its first coordinate.
Its first two moments are
\begin{equation}\label{poly:moments}
 \begin{gathered}
 \int x\,d\gamma_v(x)=-a_v,
 \qquad
 \int x^2\,d\gamma_v(x)=a_v^2+\sum_j\nu_j,
 \\
 \sum_j\nu_j=\sum_{w\ne v}|W_{vw}|^2\sum_i t_{vi}t_{wi}.
 \end{gathered}
\end{equation}

Let $N_f(x)$ and $N_g(x)$ count roots strictly greater than $x$, including
multiplicities, and define
\begin{equation}\label{poly:local-functionals}
 \xi_h(x)=N_f(x)-N_g(x),\qquad
 \ell_v=\int_0^\infty\xi_h(x)\,dx,\qquad
 m_v=\int_\R x_+\,d\gamma_v(x).
\end{equation}
Then $\xi_h=\mathbf1_{\{h<0\}}$ off a finite subset of $\R$, and
\begin{equation}\label{poly:poisson}
 \frac{1}{\pi}\arg h(x+i\eta)
   =(P_\eta*\xi_h)(x),
 \qquad P_\eta(x)=\frac{\eta}{\pi(x^2+\eta^2)},\quad\eta>0,
\end{equation}
where $\arg h\in(0,\pi)$. Moreover,
\begin{equation}\label{poly:sum-local-measures}
 \sum_{v\in Q}G_{v,Q}(z)=\frac{f'_Q(z)}{f_Q(z)},
 \qquad
 \sum_{v\in Q}\gamma_v=\sum_{f_Q(\lambda)=0}\delta_\lambda,
\end{equation}
with multiplicities on the right.

If $v$ is made deterministic of color $i$, write $h_i$ for the new quotient
relative to the same $g$, and put $R_i=h_i/h$. Its representation has the same $a_v$ and poles
$s_j$, and masses $\nu_{ij}\geq0$ which can be chosen so that
\begin{equation}\label{poly:pin-average}
 h=\sum_i t_{vi}h_i,\qquad
 \nu_j=\sum_i t_{vi}\nu_{ij}.
\end{equation}
If the penalty at $v$ is also changed to $Ka_v$, the quotient becomes
\begin{equation}\label{poly:buffered-quotient}
 h_i'=h_i+(K-1)a_v.
\end{equation}
Here the prime with a color index denotes the buffered quotient, not a
derivative. We write $\ell'_{v,i}$ for the positive-axis integral of
the spectral shift defined using $h_i'$ and the unchanged $g$.
\end{lemma}

\begin{proof}
Fix every variable except $Y_v$ at $z$ in the upper half-plane. By
\eqref{poly:deletion}, division by $g(z)\ne0$ produces a monic affine
polynomial $Y_v-\alpha(z)$. Its root cannot be in the upper half-plane by
Lemma~\ref{lem:stability}. Hence
$h(z)=z-\alpha(z)$ satisfies \eqref{poly:pick-inequality}.

The rational function $h$ has real coefficients and only real poles.
Its upper-half-plane property rules out poles of order two or higher:
the imaginary part of a real leading polar term of such order has both
signs along rays approaching the pole from above. At an actual simple
pole the residue must be negative, as seen by approaching vertically.
The polynomial part is $z+a_v$, by comparing the two highest coefficients
of the monic polynomials $f$ and $g$. This gives
\eqref{poly:arrow-quotient}. For repeated roots of $g$, assign the total
residue to one copy and pad with zero masses at the others. The Schur
complement proves \eqref{poly:arrow-matrix} and
\eqref{poly:stieltjes}, first away from the poles and then as identities
of rational functions. The first two moments follow from the first
diagonal entries of $B$ and $B^2$. Expanding $f/g$ at infinity, the
coefficient of $1/z$ is minus the sum of the length-two cycle weights
through $v$, giving the final identity in \eqref{poly:moments}.

For real $x$ away from roots and poles, the Schur complement of $xI-S$
in $xI-B$ is $h(x)$. Inertia is additive under this congruence, so the
number of negative eigenvalues of $xI-B$ exceeds that of $xI-S$ by
exactly $\mathbf1_{\{h(x)<0\}}$. This proves the assertion about
$\xi_h$ and in particular $0\leq\xi_h\leq1$ almost everywhere.

Factor $f$ and $g$ over their real roots. On the upper half-plane, take
each factor $z-\lambda$ with argument in $(0,\pi)$. The sum of the
numerator arguments minus the denominator arguments equals $\arg h$:
their difference is a continuous multiple of $2\pi$ and tends to zero
at $z=iR$ as $R\to\infty$. The argument of $x+i\eta-\lambda$, divided
by $\pi$, is the Poisson integral of
$\mathbf1_{(-\infty,\lambda)}$. Summing and subtracting gives
\eqref{poly:poisson}. Finally the chain rule and
\eqref{poly:deletion} give the first identity in
\eqref{poly:sum-local-measures}. The second follows by comparing the
finite Stieltjes transforms on the two sides.

Pinning changes only the probability at $v$, so $g$ is unchanged and
$f$ is the weighted sum of its pinned versions. This gives the first
identity in \eqref{poly:pin-average}. Each pinned polynomial is stable
by Lemma~\ref{lem:stability}, and thus has the same type of quotient
representation. Comparing residues at each distinct pole gives the
second identity; use the same placement on repeated copies for every
choice. Finally changing the $v$-penalty adds $(K-1)a_v g$ to the
polynomial, proving \eqref{poly:buffered-quotient}.
\end{proof}

The two local quantities have different roles. Define the total
positive-root mass, with roots counted according to multiplicity, by
\begin{equation}\label{poly:total-cost}
 \mathcal E(Q)=\sum_{f_Q(\lambda)=0}\lambda_+.
\end{equation}
The sum-of-measures identity and the definition of the spectral shift give
\begin{equation}\label{poly:cost-response}
 \mathcal E(Q)=\sum_{v\in Q}m_v,\qquad
 \ell_v=\mathcal E(Q)-\mathcal E(Q\setminus\{v\}).
\end{equation}
Thus $m_v$ distributes the total cost among the vertices, whereas
$\ell_v$ is the cost of inserting $v$. A buffered pin at $v$ leaves
the deleted polynomial unchanged, so its change of total cost is
exactly $\ell_{v,i}'-\ell_v$. These are costs of the expected
characteristic polynomial; in general they are not expectations of
the positive-part trace of the randomly paved matrix. They become
the ordinary spectral costs when all colors are deterministic.

\subsection{The initial largest-root bound}

The following estimate is obtained in
\cite[Lemma~24 and proof of Theorem~25]{RS21}. We include a derivation
from the MSS covariance bound, explaining both the positive-defect
completion and the deletion of the added Gram columns.

\begin{lemma}[Initial largest-root bound]\label{lem:mss-paving}
Let $T$ be a Hermitian contraction on a nonempty finite index set $Q$,
with zero diagonal. In \eqref{poly:definition}, take $W=T$, $a_v=0$,
and $t_{vi}=1/r$. Then
\begin{equation}\label{poly:mss-bound}
 \max\operatorname{root}(f_Q)\leq2\sqrt{2/r}+2/r.
\end{equation}
In particular, giving every vertex the same penalty $b$ places all
roots strictly below zero if $2\sqrt{2/r}+2/r<b$.
\end{lemma}

\begin{proof}
Set $D_0=I+T$. Then $0\leq D_0\leq2I$ and $(D_0)_{vv}=1$.
Independently form the columns
\[
 u_v=D_0^{1/2}e_v\otimes e_{\mathrm{color}(v)}
       \quad\text{in }\C^Q\otimes\C^r.
\]
They have squared norm one, their Gram matrix is $I+T_{\mathrm{pav}}$,
and, with $d=2/r$,
\[
 \sum_{v\in Q}\E u_vu_v^*=(D_0/r)\otimes I_r\leq dI.
\]
Write the positive defect as a finite sum of deterministic outer products.
Splitting its positive eigenvalues into finitely many summands
at most one makes each added column have squared norm at most one.
For all original and added columns together, their covariances $A_j$
satisfy $\sum_j A_j=dI$ and $\Tr A_j\leq1$.

Theorem~4.1 of \cite{MSS15} identifies the expected characteristic
polynomial of a sum of independent, finitely supported rank-one outer
products with the mixed characteristic polynomial of the covariances.
Theorem~5.1 of \cite{MSS15} gives largest root at most
$(1+\sqrt{\epsilon})^2$ for positive semidefinite covariances summing
to $I$ with individual traces at most $\epsilon$. Apply these results
to $A_j/d$, with $\epsilon=1/d$, and rescale the spectral variable.
The expected characteristic polynomial of the enlarged sum has largest
root at most
\begin{equation}\label{poly:enlarged-root}
 d(1+\sqrt{1/d})^2=1+2\sqrt d+d.
\end{equation}

We justify carefully the removal of the added columns. For independent,
finitely supported columns $u_1,\ldots,u_m$ in a finite-dimensional space, let
$A_j=\E u_ju_j^*$ and let $U$ be the matrix of columns. Then
\begin{equation}\label{poly:gram-stability}
 F(Y):=\E\det\bigl(\diag(Y)-U^*U\bigr)
 =\left.
   \left(\prod_{j=1}^m(Y_j-\partial_{s_j})\right)
   \det\left(I+\sum_{j=1}^m s_j A_j\right)
   \right|_{s=0}.
\end{equation}
Indeed, the coefficient indexed by a subset $J$ of columns on either
side is $(-1)^{|J|}$ times the expected principal Gram determinant
$\E\det(U_J^*U_J)$. This follows first for deterministic rank-one
matrices by determinant expansion, and then for the mean outer products
by independence and multilinearity in the distinct matrices indexed
by $J$.

The determinant on the right of \eqref{poly:gram-stability} is stable
in the variables $s_j$. If it had a null vector with all
$\operatorname{Im}s_j>0$, taking the imaginary part of its quadratic
form would force the vector into $\ker A_j$ for every $j$; the identity
term then forces the vector to be zero. If a polynomial $p$ is stable,
fixing all variables except $s_j$ at upper-half-plane arguments gives
\[
 \frac{\partial_{s_j}p}{p}=
 \sum_k\frac{1}{s_j-\lambda_k},\qquad
 \operatorname{Im}\lambda_k\leq0.
\]
The sum is zero when the degree is zero, and otherwise has negative
imaginary part. Thus a new upper-half-plane variable $Y_j$ cannot equal
this logarithmic derivative, and
$(Y_j-\partial_{s_j})p$ is stable. Repeating this observation, and then
specializing the $s_j$ to zero by the zero-free limit theorem, proves
stability of $F$. The limit is not the zero polynomial, since its
leading monomial $\prod_jY_j$ has coefficient one. The same proof
works for every principal set of columns.

Differentiating $F$ in a diagonal variable deletes the corresponding
column from the expected Gram polynomial. The monic affine argument
of Lemma~\ref{lem:local-data} therefore applies to the quotient of its
univariate specialization by the deleted polynomial. It has
imaginary part at least $\operatorname{Im}z$, and its polynomial part
is $z-\E\norm{u_j}^2$. The residue argument in that Lemma realizes
the full polynomial as the characteristic polynomial of a Hermitian
arrow matrix whose principal compression has the deleted polynomial
as characteristic polynomial. The largest eigenvalue of a principal
compression cannot exceed the largest eigenvalue of the whole matrix.
Thus deleting columns cannot increase the largest root.

If the ambient column space has dimension $q$, the exact identity
\[
 z^q\E\det(zI_m-U^*U)
   =z^m\E\det(zI_q-UU^*)
\]
shows that the enlarged Gram polynomial and the enlarged sum polynomial
differ only by zero roots. Since the bound
\eqref{poly:enlarged-root} is positive, it also bounds the largest
root of the enlarged Gram polynomial. Delete every added column using
the preceding paragraph. The original Gram polynomial is
$\E\det(zI-(I+T_{\mathrm{pav}}))=f_Q(z-1)$.
Shifting by minus one proves \eqref{poly:mss-bound}. A common penalty
$b$ replaces $f_Q(z)$ by $f_Q(z+b)$, giving the final assertion.
\end{proof}

\subsection{Residue tails and exceedance lengths}

The rest of this section concerns arbitrary arrow quotients, not just
those arising from a coloring. Let $s_j\in\R$, $\nu_j\geq0$, and
\begin{equation}\label{poly:general-arrow}
 h(z)=z+a-\sum_j\frac{\nu_j}{z-s_j},\qquad a>0.
\end{equation}
Let $\gamma$ be the first-coordinate spectral probability measure of
the corresponding arrow matrix, and put
\begin{equation}\label{poly:length-definition}
 m=\int x_+\,d\gamma(x),\qquad
 L(a,\nu)=\left|\left\{\lambda>0:
       \sum_j\frac{\nu_j}{\lambda-s_j}>\lambda+a\right\}\right|.
\end{equation}
Here and below $|\cdot|$ denotes Lebesgue measure when applied to a set
of real numbers. By Lemma~\ref{lem:local-data}, this length is the
positive-axis integral of the spectral shift for $h$.

\begin{lemma}[Control of the residues]\label{poly:tail-control}
Suppose $a=1$ in \eqref{poly:general-arrow}, and set
$H=\int|x|\,d\gamma(x)=1+2m$. There is a universal $C_0$ such that
\begin{align}
 \sum_{|s_j|\leq60H}\nu_j&\leq C_0H^2,
       \label{poly:near-mass}\\
 \sum_{s_j\geq20H}\frac{\nu_j}{s_j}&\leq C_0m,
       \label{poly:positive-tail}\\
 \sum_{s_j\leq-20H}\frac{\nu_j}{|s_j|}&\leq C_0H.
       \label{poly:negative-tail}
\end{align}
For example, $C_0=6000$ suffices.
\end{lemma}

\begin{proof}
The first spectral moment is $-1$, so $H=1+2m\geq1$.
By Markov's inequality,
\[
 \gamma([-4H,4H])\geq3/4.
\]
Writing $G=1/h$, we obtain
\[
 -\operatorname{Im}G(i10H)
 =\int\frac{10H}{x^2+100H^2}\,d\gamma(x)
 \geq\frac{15}{232H}.
\]
Since $\operatorname{Im}(1/G)\leq1/(-\operatorname{Im}G)$,
\[
 10H+\sum_j\frac{10H\nu_j}{s_j^2+100H^2}
 =\operatorname{Im}h(i10H)\leq\frac{232}{15}H<16H.
\]
For $|s_j|\leq60H$, the denominator is at most $3700H^2$.
It follows that the left side of \eqref{poly:near-mass} is at most
$2220H^2$.

Let $e$ be the first coordinate vector of the arrow matrix $B$, and let
$e_j$ be the other coordinate vectors, so that
$(B-s_j)e_j$ is a scalar multiple of $e$ of modulus $\sqrt{\nu_j}$.
Fix $\lambda\geq4H$ and $s_j\geq2\lambda$. Spectral calculus on
$(-\infty,\lambda]$, which is separated from $s_j$, and
$\ip{e_j}{e}=0$ give
\begin{align*}
 \left|\ip{e_j}{\mathbf1_{(\lambda,\infty)}(B)e}\right|
 &=\left|\ip{e_j}{\mathbf1_{(-\infty,\lambda]}(B)e}\right|\\
 &=\sqrt{\nu_j}\int_{x\leq\lambda}\frac{d\gamma(x)}{s_j-x}
 \geq\frac{\sqrt{\nu_j}}{2s_j}.
\end{align*}
The last inequality uses
$\gamma([-\lambda,\lambda])\geq3/4$ and
$s_j+\lambda\leq3s_j/2$. Bessel's inequality now yields
\[
 \sum_{s_j\geq2\lambda}\frac{\nu_j}{s_j^2}
 \leq4\gamma((\lambda,\infty)).
\]
Integrating for $\lambda\geq4H$ gives
\[
 \sum_{s_j\geq8H}\frac{\nu_j}{s_j^2}
                    \left(\frac{s_j}{2}-4H\right)
 \leq4\int_{4H}^\infty\gamma((\lambda,\infty))\,d\lambda
 \leq4m.
\]
For $s_j\geq20H$ the summand on the left is at least
$3\nu_j/(10s_j)$. Thus \eqref{poly:positive-tail} holds with constant
$40/3$. The same positive-tail calculation applied to $-B$ has first
absolute moment $H$ and positive first moment
$\int(-x)_+\,d\gamma(x)=1+m\leq H$. It gives
\eqref{poly:negative-tail} with the same constant. These calculations
include zero residues and repeated poles.
\end{proof}

The next lemma is Boole's identity for finite positive atomic measures
\cite{Boole1857}; see \cite[Proposition~2.6 and Appendix~A]{AW15}.
For this one-sided normalization, see
\cite[Section~1, equations~(1.1), (1.3), and~(1.5)]{PSZ09}.
Its elementary root-sum proof will also fix the sign and endpoint
conventions used here.

\begin{lemma}[Boole length identity]\label{poly:boole}
For finitely many real poles $s_j$, nonnegative weights $w_j$, and $c>0$,
\begin{equation}\label{poly:boole-identities}
 \left|\left\{\lambda\in\R:
       \sum_j\frac{w_j}{\lambda-s_j}>c\right\}\right|
 =\left|\left\{\lambda\in\R:
       \sum_j\frac{w_j}{\lambda-s_j}<-c\right\}\right|
 =\frac{\sum_jw_j}{c}.
\end{equation}
If every pole carrying positive weight is strictly positive, the first
exceedance set lies in $(0,\infty)$.
\end{lemma}

\begin{proof}
Merge coincident poles and omit zero weights. If there are no remaining
poles the assertion is immediate. Otherwise list them in increasing
order. The rational function is strictly decreasing between poles.
Immediately to the right of each pole it tends to $+\infty$; before
the next pole it tends to $-\infty$, and after the last pole it tends
to zero. Thus the positive-level exceedance consists of one interval
from each pole to the ensuing root of the level equation. Clearing
denominators in
$\sum_jw_j/(\lambda-s_j)=c$, comparison of the two highest
coefficients shows that the sum of the level roots exceeds the sum
of the poles by $(\sum_jw_j)/c$. This is precisely the total interval
length. Reflection of the variable gives the negative-level identity.
With positive poles the rational function is nonpositive on
$(-\infty,0]$, proving the last assertion. Values at the finitely many
poles and endpoints do not affect these lengths.
\end{proof}

\subsection{The buffered choice}

The buffer makes it possible to retain coefficient one on the existing
cost $\ell$. When $m$ is large, we bound the whole child cost by a multiple
of $m$; when $m$ is small, we reserve half of the existing length and
control the remaining exceedance by the residue tails.

\begin{lemma}[Buffered pin inequality]\label{lem:buffered-pin}
There are universal constants $K\geq1$ and $C>0$ with the following
property. In \eqref{poly:general-arrow}, let a finite family of choices
have masses $\nu_{ij}\geq0$ and probabilities $t_i\geq0$ such that
$\sum_i t_i=1$ and $\sum_i t_i\nu_{ij}=\nu_j$. Keep the poles $s_j$
unchanged, and define
\[
 h_i'(z)=z+Ka-\sum_j\frac{\nu_{ij}}{z-s_j},
 \qquad \ell_i'=L(Ka,\nu_i),\qquad \ell=L(a,\nu).
\]
Then
\begin{equation}\label{poly:buffered-pin-bound}
 \min_{i:t_i>0}\ell_i'\leq\ell+Cm.
\end{equation}
The constants have no dependence on the number or positions of the
poles, the masses, the probabilities, or the positive penalty $a$.
\end{lemma}

\begin{proof}
Dividing the arrow matrix by $a$ sends $s_j$ to $s_j/a$, the masses to
$\nu_j/a^2$, and both lengths and $m$ to their values divided by $a$.
It therefore suffices to prove the assertion for $a=1$. Write
$L=L(1,\nu)$ and $H=1+2m$, and fix $C_0$ from
Lemma~\ref{poly:tail-control}. Only positive-probability choices are
used below.

Suppose first that $m\geq1$. The means of the two nonnegative
statistics
\[
 A_i=\sum_{|s_j|\leq20H}\nu_{ij},\qquad
 B_i=\sum_{|s_j|>20H}\frac{\nu_{ij}}{|s_j|}
\]
are at most $C_0H^2$ and $C_0(m+H)\leq2C_0H$, respectively.
Markov's inequality and a union bound give a choice satisfying
\begin{equation}\label{poly:large-m-selection}
 A_i\leq C_2H^2,\qquad B_i\leq C_2H,
 \qquad C_2=8C_0.
\end{equation}
For $\lambda>DH$ and $D\geq40$, all denominators from poles
$|s_j|\leq20H$ are positive and their total contribution is at most
$C_2H^2/(\lambda-20H)\leq C_2H$. The contribution from poles
$s_j<-20H$ is at most $C_2H$ as well. Put $C_3=2C_2$.
For the remaining positive poles the identity
\begin{equation}\label{poly:large-m-transform}
 \sum_{s_j>20H}\frac{\nu_{ij}}{\lambda-s_j}
 = (\lambda+H)\sum_{s_j>20H}
       \frac{\nu_{ij}}{(s_j+H)(\lambda-s_j)}
       -\sum_{s_j>20H}\frac{\nu_{ij}}{s_j+H}
\end{equation}
holds. Choose $D\geq\max\{40,2C_3+1\}$. For any $K\geq1$, if
the entire rational sum exceeds $\lambda+K$, then
\eqref{poly:large-m-transform} forces
\[
 \sum_{s_j>20H}\frac{\nu_{ij}}{(s_j+H)(\lambda-s_j)}>\frac12.
\]
Indeed the necessary lower bound for this sum is at least
$(\lambda-C_3H)/(\lambda+H)\geq1/2$. By
Lemma~\ref{poly:boole} and \eqref{poly:large-m-selection}, its
exceedance length is at most
\[
 2\sum_{s_j>20H}\frac{\nu_{ij}}{s_j+H}\leq2C_2H.
\]
Including the initial
interval $(0,DH)$, we obtain
\begin{equation}\label{poly:large-m-conclusion}
 L(K,\nu_i)\leq(D+2C_2)H\leq3(D+2C_2)m.
\end{equation}
This proves the required bound in this case, with a universal constant,
for every $K\geq1$.

It remains to treat $0\leq m<1$, so $1\leq H\leq3$. Set
\[
 P=\sum_{s_j>0}\frac{\nu_j}{1+s_j},\qquad
 M_-(d)=\sum_{s_j\leq0}\frac{\nu_j}{d+|s_j|}\quad(d>0),
\]
using the right limit $M_-(0+)$ when necessary. Splitting the positive
poles at $60$ and using Lemma~\ref{poly:tail-control} gives
\begin{equation}\label{poly:P-bound}
 P\leq9C_0+C_0m\leq C_4,
 \qquad C_4=10C_0.
\end{equation}
Here $60\leq60H$ controls the near part and $60\geq20H$ controls
the tail. For $\lambda>0$ away from the poles, the parent exceedance
condition is equivalent to
\begin{equation}\label{poly:parent-exceedance}
 R(\lambda):=\sum_{s_j>0}
 \frac{\nu_j}{(1+s_j)(\lambda-s_j)}
 >1+\frac{P-M_-(\lambda)}{1+\lambda}.
\end{equation}
The threshold on the right is at most $1+P$. All poles of $R$ with
positive weight are positive. Lemma~\ref{poly:boole} therefore gives
\begin{equation}\label{poly:P-controlled-by-L}
 L\geq\frac{P}{1+P},\qquad
 P\leq C_5L,\qquad C_5=1+C_4.
\end{equation}

We next prove the uniform bound
\begin{equation}\label{poly:negative-control}
 M_-(L/2)\leq C_6,
\end{equation}
with the right-limit interpretation when $L=0$. If $L\geq1$, split
the negative poles at absolute value $60H$. The near part is at most
$2C_0H^2\leq18C_0$, while the far part is at most $C_0H\leq3C_0$.
Thus $21C_0$ suffices in this case. If $L=0$, then
\eqref{poly:P-controlled-by-L} implies $P=0$. If $M_-(0+)>1$,
including an infinite right limit, then
$M_-(\lambda)>1+\lambda$ on a nonempty interval immediately to the
right of zero. With no positive-pole contribution this would be an
interval of parent exceedance, contradicting $L=0$. Hence
$M_-(0+)\leq1$.

Finally suppose $0<L<1$. If
$M_-(2L)>7C_5+9$, then for $0<\lambda<2L$ we have
$P\leq C_5$, $M_-(\lambda)\geq M_-(2L)$, and $1+\lambda\leq3$.
Consequently the right side of \eqref{poly:parent-exceedance} is
strictly less than $-2(C_5+1)$. By the negative-level identity in
Lemma~\ref{poly:boole},
\[
 \left|\{\lambda\in\R:R(\lambda)<-2(C_5+1)\}\right|
 =\frac{P}{2(C_5+1)}<\frac L2.
\]
Outside this exceptional set, the parent exceedance holds throughout
$(0,2L)$, apart from finitely many poles. Its length is therefore
greater than $3L/2$, contradicting the definition of $L$. It follows
that $M_-(2L)\leq7C_5+9$. Termwise comparison of denominators gives
$M_-(L/2)\leq4M_-(2L)$. Thus
\eqref{poly:negative-control} holds in all cases with, for instance,
\begin{equation}\label{poly:C6-choice}
 C_6=\max\{21C_0,1,28C_5+36\}.
\end{equation}

For each choice define $P_i$ and $M_{-,i}$ by replacing $\nu$ with
$\nu_i$. Averaging these nonnegative quantities, followed by Markov's
inequality and a union bound, gives a single choice $i$ with
\begin{equation}\label{poly:small-m-selection}
 P_i\leq4P,\qquad
 M_{-,i}(L/2)\leq4C_6,\qquad
 \sum_{s_j>60}\frac{\nu_{ij}}{s_j}\leq4C_0m.
\end{equation}
There is no exception when one of the means is zero: the corresponding
nonnegative statistic then vanishes at every positive-probability
choice. When $L=0$, finiteness of $M_-(0+)$ forces $\nu_j=0$ at every
pole $s_j=0$. The nonnegative averaging identities then force
$\nu_{ij}=0$ at those poles for every positive-probability choice.
The right-limit child statistics are therefore finite sums over
strictly negative poles, and their averages equal the parent statistic.
Thus the same selection
argument applies at $L=0$ and at $m=0$.

For this choice put $P_i^{(K)}=\sum_{s_j>0}\nu_{ij}/(K+s_j)$.
The child's exceedance condition is
\begin{equation}\label{poly:child-exceedance}
 \sum_{s_j>0}\frac{\nu_{ij}}{(K+s_j)(\lambda-s_j)}
 >1+\frac{P_i^{(K)}-M_{-,i}(\lambda)}{K+\lambda}.
\end{equation}
For $\lambda>L/2$ the negative statistic is at most $4C_6$ by
monotonicity. Therefore $K\geq8C_6$ makes the right side of
\eqref{poly:child-exceedance} at least $1/2$. The length of the child
exceedance beyond $L/2$ is at most $2P_i^{(K)}$ by
Lemma~\ref{poly:boole}. Split the positive poles at $60$. For
$0<s_j\leq60$,
\[
 \frac{1}{K+s_j}\leq\frac{61}{K}\frac{1}{1+s_j},
\]
whereas for $s_j>60$ it is at most $1/s_j$. Thus
\begin{equation}\label{poly:child-length-estimate}
 2P_i^{(K)}
 \leq\frac{122}{K}P_i
       +2\sum_{s_j>60}\frac{\nu_{ij}}{s_j}
 \leq\frac{488C_5}{K}L+8C_0m.
\end{equation}
Choose once and for all
\[
 K\geq\max\{1,8C_6,976C_5\}.
\]
The first term in \eqref{poly:child-length-estimate} is at most $L/2$.
Adding the possible length $L/2$ of the initial interval gives
$L(K,\nu_i)\leq L+8C_0m$. Together with
\eqref{poly:large-m-conclusion}, this proves
\eqref{poly:buffered-pin-bound} with
$C\geq\max\{8C_0,3(D+2C_2)\}$. Scaling back to arbitrary $a>0$
completes the proof.
\end{proof}

We fix the universal constants $K,C$ supplied by
Lemma~\ref{lem:buffered-pin} for the remainder of the argument.
By \eqref{poly:cost-response}, at a vertex with positive penalty the
lemma supplies a buffered pin increasing $\mathcal E(Q)$ by at most
$Cm_v$. The coefficient one on
$\ell_v$ is what cancels the old insertion cost in this difference.
The measured argument will sum these local charges while making
interactions between simultaneous pins negligible.
